% ======================================================================
%  ECCV Supplementary — Detailed Module Descriptions
%  Sections: Image and Prompt Encoding  /  Lightweight Task-Specific Decoding
% ======================================================================

\section{Detailed Module Descriptions}
\label{sec:supp_modules}

This section provides the detailed internal structures of the two
lightweight stages that wrap around the core interleaved-attention
backbone described in the main paper (Sec.~\ref{subsec:pipeline}).

% ======================================================================
\subsection{Image and Prompt Encoding}
\label{sec:supp_enc}

\noindent\textbf{Image Encoder.}
Both the reference (satellite) and query images are resized to $518{\times}518$ and normalized with ImageNet statistics. A frozen DINOv2-ViT-L/14 encoder~\cite{dinov2} processes each image independently, producing $N_p = (518/14)^2 = 1369$ patch tokens of dimension $D_\text{enc}=1024$:
\begin{equation}
  T_r, T_q \in \mathbb{R}^{N_p \times D_\text{enc}}.
\end{equation}
Five learnable register tokens~\cite{dinov2reg} are prepended to each patch sequence inside the decoder (see Sec.~3 of the main paper), bringing the per-view token count to $N_v = 5 + N_p = 1374$.

\noindent\textbf{Prompt Encoder.}
We adapt the SAM-2 prompt encoder~\cite{sam2} with a dual-dimension design: all internal layers operate at the native SAM dimension $d_\text{sam}=256$, which allows direct loading of pretrained SAM-2 weights; a pair of learnable projection layers then maps the output to the decoder dimension $D = 1024$. The three prompt modalities are handled as follows.

\emph{(i) Point prompt.}
Given a single pixel coordinate $(x, y)$ in the query image, we first shift by $+0.5$ to center the pixel and normalize to $[0,1]$:
\begin{equation}
  \bar{\mathbf{c}} = \bigl(\tfrac{x+0.5}{W},\;\tfrac{y+0.5}{H}\bigr).
\end{equation}
A random Fourier positional encoding maps the 2-D coordinate to a $d_\text{sam}$-dimensional vector:
\begin{equation}
  \text{PE}(\bar{\mathbf{c}})
    = \bigl[\sin(2\pi\,\mathbf{G}\,\bar{\mathbf{c}}'),\;
            \cos(2\pi\,\mathbf{G}\,\bar{\mathbf{c}}')\bigr]
    \in \mathbb{R}^{d_\text{sam}},
\end{equation}
where $\bar{\mathbf{c}}' = 2\bar{\mathbf{c}}-1$ and $\mathbf{G}\in\mathbb{R}^{(d_\text{sam}/2)\times 2}$ is a frozen Gaussian matrix drawn at initialization.
A learned type embedding $\mathbf{e}_\text{pos}\in\mathbb{R}^{d_\text{sam}}$ (foreground-point type) is added, and a padding token with a dedicated \emph{not-a-point} embedding is appended when no box prompt is present, yielding the sparse token(s)
\begin{equation}
  T_p^{(\text{pt})} \in \mathbb{R}^{(1{+}1)\times d_\text{sam}}\quad
  \text{(point + padding)}.
\end{equation}

\emph{(ii) Bounding-box prompt.}
A box $(x, y, w, h)$ in pixel space is converted to two corner coordinates
$\mathbf{c}_1=(x,y)$, $\mathbf{c}_2=(x{+}w,\,y{+}h)$.
Each corner is encoded identically to a point (shift, normalize, Fourier PE) and summed with a corner-type embedding ($\mathbf{e}_\text{tl}$ for the top-left corner, $\mathbf{e}_\text{br}$ for the bottom-right corner):
\begin{equation}
  T_p^{(\text{box})} =
    \bigl[\text{PE}(\bar{\mathbf{c}}_1)+\mathbf{e}_\text{tl},\;
          \text{PE}(\bar{\mathbf{c}}_2)+\mathbf{e}_\text{br}\bigr]
    \in \mathbb{R}^{2\times d_\text{sam}}.
\end{equation}

\emph{(iii) Mask prompt.}
A binary mask $M\in\{0,1\}^{H\times W}$ is first bilinearly resized to
$4 H_p \times 4 W_p = 148{\times}148$ and then processed by a lightweight
CNN that progressively down-samples while expanding channels:
\begin{equation}
  M
  \;\xrightarrow[\text{stride-2 Conv}]{1{\to}4}
  \;\xrightarrow{\text{LN + GELU}}
  \;\xrightarrow[\text{stride-2 Conv}]{4{\to}16}
  \;\xrightarrow{\text{LN + GELU}}
  \;\xrightarrow[\text{1\!\times\!1 Conv}]{16{\to}d_\text{sam}}
  E_d \in \mathbb{R}^{d_\text{sam}\times H_p\times W_p}.
\end{equation}
The two stride-2 convolutions reduce the spatial size from $148^2$ to $74^2$ to $37^2=H_p\!\times\!W_p$, matching the patch grid.
When no mask prompt is given, a learned no-mask embedding $\mathbf{e}_\emptyset\in\mathbb{R}^{d_\text{sam}}$ is broadcast to the full $H_p\!\times\!W_p$ spatial grid.

\noindent\textbf{Dimension Projection.}
All SAM-internal representations share the dimension $d_\text{sam}=256$.
To interface with the $D=1024$-dimensional backbone decoder, two learnable projection layers—the only randomly initialized parts of the prompt encoder—are applied:
\begin{align}
  T_p &= \text{LN}\!\bigl(\mathbf{W}_s\, T_p^{(*)}\bigr)
       \in \mathbb{R}^{N_e\times D},
       &\mathbf{W}_s &\in\mathbb{R}^{D\times d_\text{sam}},
       \label{eq:sparse_proj}\\
  E_d &= \text{LN2d}\!\bigl(\text{Conv}_{1\times1}(E_d^{(*)})\bigr)
       \in \mathbb{R}^{D\times H_p\times W_p},
       \label{eq:dense_proj}
\end{align}
where LN denotes Layer Normalization and LN2d its spatial variant.

\noindent\textbf{Learnable Task Tokens.}
In addition to the prompt tokens, we introduce a set of learnable task tokens $T_l\in\mathbb{R}^{N_l\times D}$, initialized from $\mathcal{N}(0, 0.02^2)$ and jointly optimized during training. These tokens are partitioned into two functional groups:
\begin{equation}
  T_l = [T_\text{obj},\; T_\text{pos}],\quad
  T_\text{obj}\in\mathbb{R}^{N_\text{obj}\times D},\quad
  T_\text{pos}\in\mathbb{R}^{N_\text{pos}\times D},
\end{equation}
with $N_\text{obj}=1$ and $N_\text{pos}=1$ in our default configuration.
$T_\text{obj}$ drives detection and segmentation outputs, while $T_\text{pos}$ serves the position heatmap prediction. All learnable tokens are concatenated to the reference stream during local attention (see the main paper); their RoPE positions are set to the spatial origin $(0,0)$.


% ======================================================================
\subsection{Lightweight Task-Specific Decoding}
\label{sec:supp_dec}

After the interleaved local--global attention backbone, we obtain updated
reference features $\hat{T}_r\in\mathbb{R}^{N_p\times D_\text{out}}$,
query features $\hat{T}_q\in\mathbb{R}^{N_p\times D_\text{out}}$,
and task tokens $\hat{T}_l\in\mathbb{R}^{N_l\times D_\text{out}}$,
where $D_\text{out}=2D=2048$ is obtained by concatenating the outputs of
the last two decoder layers and projecting through a linear layer.
We partition $\hat{T}_l = [\hat{T}_\text{obj},\;\hat{T}_\text{pos}]$ and
feed them into four lightweight heads.

% ------ Detection head ------
\subsubsection{Detection Head $\mathcal{D}_\text{det}$}
We adopt a DETR-style~\cite{DETR} detection head.
Box regression uses a 3-layer MLP with ReLU activations:
\begin{equation}
  \hat{b} = \sigma\!\bigl(\text{MLP}_3(\hat{T}_\text{obj})\bigr)
           \in [0,1]^{N_\text{obj}\times 4},
  \quad
  \text{MLP}_3: D_\text{out}\to D_\text{out}\to D_\text{out}\to 4,
\end{equation}
where $\sigma$ is the sigmoid function and the four outputs encode the
normalized center coordinates and dimensions $(c_x, c_y, w, h)$.
A parallel linear classifier produces a single-class objectness score:
\begin{equation}
  s = \sigma\!\bigl(\mathbf{w}_\text{cls}^\top \hat{T}_\text{obj}\bigr)
    \in [0,1]^{N_\text{obj}}.
\end{equation}
When $N_\text{obj}>1$, Hungarian matching~\cite{DETR} assigns predictions to
the ground-truth target for loss computation.

% ------ Segmentation head ------
\subsubsection{Segmentation Head $\mathcal{D}_\text{seg}$}
Following SAM~\cite{SAM}, we employ a hyper-network that
dynamically generates per-query convolution kernels from $\hat{T}_\text{obj}$.
The pipeline is:

\begin{enumerate}[leftmargin=*,topsep=2pt,itemsep=1pt]
\item \textbf{Spatial upsampling.}
      The reference features are reshaped to a 2-D grid
      $\hat{F}_r\in\mathbb{R}^{D_\text{out}\times H_p\times W_p}$
      and progressively upscaled via two transposed-convolution stages:
      \begin{equation}
        \hat{F}_r
          \;\xrightarrow[\text{stride-2 DeConv}]{D_\text{out}{\to}\tfrac{D_\text{out}}{4}}
          \;\xrightarrow{\text{LN2d + GELU}}
          \;\xrightarrow[\text{stride-2 DeConv}]{\tfrac{D_\text{out}}{4}{\to}\tfrac{D_\text{out}}{8}}
          \;\xrightarrow{\text{GELU}}
          \hat{F}_r^{\uparrow}
          \in \mathbb{R}^{\tfrac{D_\text{out}}{8}\times 4H_p\times 4W_p},
      \end{equation}
      \ie $37^2\!\to\!74^2\!\to\!148^2$.

\item \textbf{Hyper-network kernel.}
      A 3-layer MLP maps each object token to a dynamic convolution kernel:
      \begin{equation}
        \mathbf{k}_i = \text{MLP}_3(\hat{T}_\text{obj}^{(i)})
                       \in \mathbb{R}^{D_\text{out}/8},
        \quad i=1,\dots,N_\text{obj}.
      \end{equation}

\item \textbf{Mask prediction.}
      A dot product between the kernel and the upscaled reference
      features produces a single-channel mask per query:
      \begin{equation}
        \hat{M}_i = \mathbf{k}_i^\top\,
                    \hat{F}_r^{\uparrow}[\cdot]
                    \in \mathbb{R}^{4H_p\times 4W_p},
      \end{equation}
      which is bilinearly interpolated to the full image resolution
      $518{\times}518$.

\item \textbf{IoU prediction.}
      A separate 3-layer MLP with sigmoid output estimates the mask
      quality for each query:
      \begin{equation}
        \widehat{\text{IoU}}_i
          = \sigma\!\bigl(\text{MLP}_3^{(\text{iou})}(\hat{T}_\text{obj}^{(i)})\bigr)
          \in [0,1].
      \end{equation}
\end{enumerate}

% ------ Position head ------
\subsubsection{Position Head $\mathcal{D}_\text{pos}$}
\label{sec:supp_pos_head}
The position head localizes the query camera on the reference image.
It follows a Mask2Former-inspired~\cite{mask2former} dynamic convolution design:

\begin{enumerate}[leftmargin=*,topsep=2pt,itemsep=1pt]
\item \textbf{Query aggregation.}
      If multiple position tokens exist ($N_\text{pos}>1$), they are merged
      via learned attention weights:
      \begin{equation}
        \hat{T}_\text{pos}^{*}
          = \sum_{j=1}^{N_\text{pos}} \alpha_j\,\hat{T}_\text{pos}^{(j)},
        \qquad
        \alpha = \text{softmax}\!\bigl(\mathbf{w}_\text{agg}^\top
                 \hat{T}_\text{pos}\bigr)
                 \in \mathbb{R}^{N_\text{pos}},
      \end{equation}
      where $\mathbf{w}_\text{agg}\in\mathbb{R}^{D_\text{out}}$ is a
      learnable weight vector.

\item \textbf{Spatial heatmap.}
      The aggregated query is dot-producted against the reference
      patch features (without upsampling) to produce a low-resolution
      logit map:
      \begin{equation}
        H(i,j) = \langle \hat{T}_\text{pos}^{*},\;
                         \hat{T}_r(i,j)\rangle,
        \qquad
        H\in\mathbb{R}^{H_p\times W_p}.
        \label{eq:heatmap_logit}
      \end{equation}

\item \textbf{Probability map.}
      $H$ is bilinearly upsampled to $518{\times}518$ and passed through
      a spatial softmax to obtain a probability distribution:
      \begin{equation}
        P = \text{softmax}\!\bigl(\text{upsample}(H)\bigr)
          \in \mathbb{R}^{518\times 518}.
      \end{equation}

\item \textbf{Coordinate extraction.}
      A differentiable soft-argmax extracts the expected 2-D position:
      \begin{equation}
        \hat{x} = \sum_{i,j} P(i,j)\cdot u_j,
        \qquad
        \hat{y} = \sum_{i,j} P(i,j)\cdot v_i,
      \end{equation}
      where $u_j=j/(W{-}1)$ and $v_i=i/(H{-}1)$ are normalized grid
      coordinates.
\end{enumerate}

\noindent\textbf{Training supervision.}
During training, the low-resolution logit map $H$ (Eq.~\ref{eq:heatmap_logit}) is supervised with a CornerNet-style~\cite{cornernet} pixel-wise focal loss (see Sec.~3.5 of the main paper), rather than applying MSE only to the extracted position. This provides dense spatial gradients and more stable optimization.

% ------ Camera head ------
\subsubsection{Camera Orientation Head $\mathcal{D}_\text{cam}$}
\label{sec:supp_cam_head}
Following~$\pi^3$~\cite{wang2025pi}, the camera head estimates the
relative orientation between the query and reference views.
It is a self-contained module composed of a per-view transformer
decoder and a pose regression network:

\begin{enumerate}[leftmargin=*,topsep=2pt,itemsep=1pt]
\item \textbf{Per-view Transformer Decoder.}
      A weight-shared transformer decoder processes the patch features
      of each view independently. It consists of a linear projection
      $D_\text{out}\to D_\text{dec}$ ($2048\to 1024$), followed by
      $L_\text{cam}=5$ standard transformer blocks with
      $\text{dim}=1024$, $16$ attention heads, MLP ratio $4$, and
      RoPE ($f_\text{rope}=100$); a final linear layer maps to
      $D_\text{cam}=512$:
      \begin{equation}
        h_v = \text{Linear}_{1024\to 512}\!\bigl(
              \text{Block}_{1..5}(\text{Linear}_{2048\to 1024}(\hat{T}_v))\bigr)
              \in\mathbb{R}^{N_p\times 512},
              \quad v\in\{q,r\}.
      \end{equation}

\item \textbf{Pose Regression (per view).}
      The $N_p$ token vectors are processed by a stack of two $1{\times}1$
      \emph{ResConv} blocks (each block: three linear layers with ReLU
      and a residual skip), followed by adaptive average pooling to a
      single vector, a 2-layer MLP ($512\to 512\to 512$, with ReLU),
      and two parallel linear heads:
      \begin{align}
        \mathbf{t}_v &= \mathbf{W}_t\,\bar{h}_v \in\mathbb{R}^{3}, \\
        \hat{\mathbf{R}}_v &= \text{SVD-ortho}\bigl(\mathbf{W}_r\,\bar{h}_v\bigr)
                              \in\text{SO}(3),
      \end{align}
      where $\bar{h}_v$ is the pooled feature, $\mathbf{W}_t\in\mathbb{R}^{3\times 512}$ and $\mathbf{W}_r\in\mathbb{R}^{9\times 512}$. The $9$-dimensional rotation output is reshaped to $3{\times}3$ and projected onto SO(3) via SVD orthogonalization~\cite{svdrotation}: $U\Sigma V^\top\!\to\!V\,\text{diag}(1,1,\det(VU^\top))\,U^\top$.

\item \textbf{Relative pose.}
      The absolute poses $(\hat{\mathbf{R}}_q,\mathbf{t}_q)$ and
      $(\hat{\mathbf{R}}_r,\mathbf{t}_r)$ are composed into a
      relative transformation $\mathbf{T}_{q\to r}=\mathbf{T}_r^{-1}\mathbf{T}_q\in\text{SE}(3)$, from which Euler angles
      $(\text{yaw},\text{pitch},\text{roll})$ are extracted using the
      ZYX convention.
\end{enumerate}

% ------ Deep Supervision ------
\subsubsection{Deep Supervision Heads}
\label{sec:supp_deep}

To improve gradient flow through the 36-layer decoder, we attach
auxiliary heads at three intermediate decoder stages
$\mathcal{S}\!=\!\{4,11,17\}$ (where one stage = one local block +
one global block). At each supervised stage $s$, a dedicated linear
projection $\text{Proj}_s: D\!\to\!D_\text{out}$ maps the single-layer
($D{=}1024$) intermediate representations to the output dimension
$D_\text{out}{=}2048$. The projected tokens are then fed into
\textbf{separate, non-shared} instances of $\mathcal{D}_\text{det}$
and $\mathcal{D}_\text{seg}$. For the early and middle stages ($s{\in}\{4,11\}$), the shared final-layer $\mathcal{D}_\text{pos}$
and $\mathcal{D}_\text{cam}$ are additionally applied to provide
heatmap and rotation supervision. The final-stage ($s{=}17$) head is used only for detection and segmentation as its output is close to the
final prediction.

Each intermediate loss is weighted by a stage factor
$w_s\in\{0.1, 0.3, 0.6\}$ (for stages $4, 11, 17$ respectively)
before being aggregated with the corresponding per-task loss weights.

